\documentclass[10pt,a4paper]{amsart}
\usepackage[margin=19mm]{geometry}
\usepackage{amsmath,amssymb,mathtools}
\usepackage[scr=rsfs,cal=euler]{mathalfa}
\usepackage{xfrac}
\usepackage[unicode,hidelinks,bookmarks=false]{hyperref}
\newcommand{\ie}{i.e.\ }
\newcommand{\cf}{cf.\ }
\newcommand{\resp}{respectively\ }
\newcommand{\Subsection}{\subsection}
\providecommand{\nobreakdash}{\nobreak\hbox{-}}
\setlength{\emergencystretch}{2em}
\multlinegap0pt
\usepackage[shortlabels]{enumitem}
\usepackage{amssymb}
\newtheorem{lemma}{Lemma}
\newcommand{\Z}{\mathbb{Z}}
\title{A geometric proof of Lagrange's theorem for continued fractions}
\author{Anton Lukyanenko, Joseph Vandehey}
\date{Source: arXiv:2603.12425v1 (CC BY 4.0)}
\begin{document}
\maketitle

\noindent\emph{Source and licence.} A geometric proof of Lagrange's theorem for continued fractions. Version arXiv:2603.12425v1 (CC BY 4.0), distributed by arXiv under the Creative Commons Attribution 4.0 International licence, http://creativecommons.org/licenses/by/4.0/, as recorded in the arXiv licence metadata for this exact version and shown on the abstract page https://arxiv.org/abs/2603.12425v1. Full licence notice: \texttt{licenses/arxiv-2603.12425-CC-BY-4.0.txt}.\par\smallskip

\noindent\textit{Editorial scope.} Counted unit: the article's lemma lemma (source0.tex lines 669--671). The article's complete original proof of it is delivered with it (source0.tex lines 673--682, one \texttt{proof} environment of 10 lines). No mathematics is written, repaired or re-derived: the statement and the proof are the article's own lines, in the article's own notation, and no retained line is substituted or re-flowed.  No further source-local context is needed: every cross-reference of every retained line resolves inside the two delivered spans.  Nothing was omitted from any retained span, so the package carries no omission record.  No retained line contains a citation, so the wrapper carries no bibliography and no bibliography entry.  No retained line contains a figure environment, an image-inclusion command or drawing code, so no image is delivered and the package declares no asset dependency.  Editorial formatting only, in addition: an AMS \texttt{amsart} preamble that restates the article's own declarations and macros used by the retained lines, namely the environments \texttt{lemma} on separate counters, together with the standard packages those lines need.  The retained environments print in this document as Lemma 1 in order of appearance, whereas the article prints the same items with the numbering of its own full document.  The complete native source of the article as distributed in the arXiv e-print for this exact version is bundled unchanged as \texttt{sources/196276/source0.tex}.
\begin{lemma} The group $\Gamma_\mathcal{H}$ is discrete; indeed, 
$\Gamma_\mathcal{H}\cup\{0\}\subset \mathbb{Z}^8 +h\mathbb{Z}^8$.
\end{lemma}

\begin{proof}
It suffices to show that $\mathbb{Z}^8 +h\mathbb{Z}^8$ is closed under multiplication. Write $a \equiv b$ if $a-b\in \Z^8$. It is straightforward to check that $h^2\equiv h$; that for each $i$ one has $e_i h \equiv h e_i$; and thus that $h\alpha \equiv \alpha h$ for any $\alpha \in \Z^8$. One then computes, for $\alpha_1+h\beta_1, \alpha_2+h\beta_2\in \Z^8 + h\Z^8$,
\begin{align*}
(\alpha_1+h\beta_1)(\alpha_2+h\beta_2)
& \equiv h\beta_1\alpha_2+\alpha_1h\beta_2+h\beta_1h\beta_2\\
&\equiv h\beta_1\alpha_2 + (h\alpha_1+\gamma_1)\beta_2+h(h\beta_1+\gamma_2)\beta_2 \\
&\equiv h (\beta_1\alpha_2+\alpha_1\beta_2+\beta_1\beta_2 +\gamma_2\beta_2),
\end{align*}
where $\gamma_1=\alpha_1h-h\alpha_1$, $\gamma_2=\beta_1 h-h\beta_1$ are both in $\Z^8$. This shows the product $(\alpha_1+h\beta_1)(\alpha_2+h\beta_2)$ is again in $\Z^8+h\Z^8$.
\end{proof}

\end{document}
